% the abstract
\begin{abstracten}
In the context of large-model self-evolution, the distribution of training tasks may be skewed and new data cannot guarantee balanced coverage, so directly applying reinforcement learning or fine-tuning degrades previously acquired capabilities. To address this, we propose a continual reinforcement learning method for agentic large language models with four components: (1) a composite continual-learning objective that combines the GRPO policy gradient with a reverse-KL anchor, experience replay, and an entropy regularizer; (2) a frozen external LLM judge that scores each trajectory on four dimensions---completion, correctness, trajectory quality, and safety---aggregated multiplicatively with safety as a binary gate; (3) a nine-bucket experience replay buffer partitioned by capability domain, with square-root-weighted capacity, independent in-bucket eviction, and distance-based bucket selection in a multidimensional capability space; and (4) dual-model task difficulty annotation with sequential bucket-by-bucket training. The method is injected into the verl reinforcement learning framework with no framework modification. Experiments are conducted on a Qwen3.5-9B base model with the ClawEval pure-text subset (195 tasks, 9 capability domains) under a train-then-evaluate-once anti-forgetting protocol, compared against a same-scale baseline without continual learning. Experimental results show that the proposed method effectively mitigates the forgetting of earlier-trained capabilities.
\end{abstracten}
